\section{Cross-check against a published calculation}
To test the model against something it was not tuned on, we compare it with arXiv:2104.14896, a published \emph{theory} study of the sodium D lines in a
wide range of magnetic fields (\texttt{src/literature.py}). That paper includes hyperfine structure (nuclear spin $3/2$). At $B=0.9\ \si{T}$ the field is
about $14$ times the hyperfine field $B_0=A_\text{hf}/\mu_B=63.3\ \si{mT}$, so the lines gather into groups, one per fine-structure Zeeman component, whose spacings
are set by the fine-structure Zeeman effect alone. Our model has no hyperfine structure, so the spacings between groups are what can be compared.

\begin{table}[h]
\centering
\begin{tabular}{lrrrr}
\toprule
Group & Spacing formula & Paper formula (GHz) & Paper numerical (GHz) & Model (GHz)\\
\midrule
D1 $\pi$ & $\tfrac43\mu_B B$ & 16.80 & 16.9 & 16.83\\
D1 $\sigma$ & $\tfrac83\mu_B B$ & 33.59 & 33.6 & 33.61\\
D2 $\pi$ & $\tfrac23\mu_B B$ & 8.40 & 8.46 & 8.42\\
D2 $\sigma$, inner pair & $2\mu_B B$ & 25.19 & 25.15 & 25.19\\
D2 $\sigma$, outer pair & $\tfrac{10}3\mu_B B$ & 41.98 & 42.07 & 42.03\\
\bottomrule
\end{tabular}
\caption{Spacings between line groups at $0.9\ \si{T}$. All model values lie within $0.5\,\%$ of the paper's.}
\end{table}

\begin{figure}[h]
\centering
\includegraphics[width=0.85\textwidth]{literature_sodium_groups.pdf}
\caption{Difference of the model from the published values for each group.}
\end{figure}

\paragraph{Limits.} This is agreement with another calculation, not with laboratory data. Some of the paper's values are read from its figures, and
the numbers above were quoted from its text (check them before citing). Hyperfine structure shifts and broadens each group by a few hundred MHz,
so about $0.5\,\%$ of a $17$ to $42\ \si{GHz}$ spacing is the expected level of agreement.
